\section{Introduction}

% The imprtance of language and what they do.

Model-driven engineering of software-intensive systems is heavily based on the definition and use of \textit{conceptual modeling languages}. Such languages prescribe how modelers are supposed to build models and how readers are meant to understand those models. A well-designed conceptual modeling language supports the efficient development of comprehensible conceptual models that are usable in different activities of the engineering lifecycle, including stakeholder communication, decision making, validation and verification, and code generation.

% Constitutents of Languages
The heart of a modeling language is an \textit{ontology} that captures a shared \textit{conceptualization} of a domain among language users \cite{Guiz.07}. The constituents of the ontology are \textit{concepts} which are used to organize the apprehension and communication of pertinent phenomena in the domain. In the language, concepts are represented through \textit{terms}, i.e., lexical, symbolic, and/or visual elements that represent concepts. UML state diagrams \cite{UML.17}, for example, are models that are based on a language containing such concepts as \textit{state}, \textit{transition}, and \textit{guard condition}, represented lexically in English through a \textit{vocabulary} of terms, such as ``state'', ``transition'', ``guard condition'', as well as designated visual elements, such as rounded rectangles, arrows, and bracketed annotations on arrows. For simplicity we will, henceforth, refer to any kind of signification of concepts (lexical, visual, or other) as \textit{terms}. 

% The design problem
%Conceptual modeling languages are artifacts, embodying a series of design decisions made by the language designers. These decisions involve, among other things, 

The selection of concepts that should be included in the language and the terms for representing those concepts, are the result of a design effort. As in any design effort, the outcome can enjoy different degrees of success as reflected by the \textit{quality} of the resulting modeling language. In recognition of this, several efforts to conceptualize language quality can be found in the literature \cite{Krog.12,NePG.12,WeWa.93}. Key focus has been the \textit{appropriateness} of a chosen vocabulary vis-a-vis the concepts of the underlying ontology it is intended to represent -- often referred to as \textit{semantic transparency} \cite{Mood.10} seen as part of the language's general \textit{comprehensibility appropriateness} \cite{Krog.12}. Wand and Weber, specifically, introduce a framework for characterizing pathologies of language designs based on how exactly the vocabulary-ontology mapping is missaligned \cite{WeWa.93}. For example, a language suffers from \textit{construct deficit} when there are concepts in the language that are not represented by corresponding terms in the vocabulary, and, reversely, \textit{construct excess}, when there are terms not obviously mapping to a concept of interest.

% The peira approach
With the dimensions of vocabulary quality theoretically defined in such a way, the natural question is how one can measure a given language design with respect to each of those dimensions. An empirical approach for doing so has been proposed by Liaskos et al. \cite{LZMK.24}. According to their framework, called \textit{Peira}, such assessments are possible via exposing experimental participants to chunks of ``raw'' unstructured domain information and asking them to classify such chunks to terms of the language vocabulary under investigation, after first exposing them to the language through a training process of specified length, depth, and intensity. The outcomes of the classification are then used to calculate a variety of metrics either based on comparison to authoritative classifications, i.e., classifications claimed to be correct by the vocabulary designers themselves, or, ignoring authoritative input, based on measures of agreement among the participants. The authors go on to propose specific formulae by which the resulting classification datasets can be used to detect the presence of language issues such as the ones discussed above -- construct deficit, construct excess, etc. 
%\red{According to the authors, the process is meant to streamline language design evaluation allowing evaluators to utilize a common set of metrics, enable comparisons between designs under certain assumptions, and avoid the subjectivity and biases lurking in alternative procedures, such as expert evaluation of complete models.}

The process is, however, not without cost given that a number of human participants need to be recruited, trained to the language, and engage in classification tasks. With the advent of Large Language Models (LLMs) \cite{LLM.survey} it is, hence, pertinent to ask whether LLMs can be used as a substitute for human participants. Specifically, instead of training humans and asking them to perform classification tasks, present the training material to an LLM and subsequently ask the LLM to perform the classification tasks instead -- all within one or a sequence of prompts. For such a strategy to be applicable at least two issues need to be clarified: (a) whether LLM performance is comparable with human performance to allow for such substitution and (b) whether the choice of LLM affects such performance and substitutability. 

In this paper, we report on a preliminary empirical study in which the Peira framework is applied with both humans and LLMs, as a first attempt to address the above questions
% via a direct comparison between the two kinds of participants. 
An example requirements modeling language that was recently introduced in the literature is adopted as the object of investigation. The Peira method is applied using on-line participants who complete a sequence of tasks in a survey-style on-line instrument, as prescribed by the method. 
%The latter consists of video presentations of the language in question as well as other instructions followed by classification tasks based on two distinct cases. 
The instrument is then converted to an LLM prompt that conveys, to the extend feasible,
% via verbatim transcription of the training video and adaptation of question and response inventory format to prompt design conventions, 
the exact same information as the on-line instrument. The LLM prompt is offered to a set of LLMs with different sizes. The resulting data sets are compared with respect to classification performance, i.e., the levels of precision and recall attained by human vs. LLM ratings when compared to authoritative ratings. We find that LLMs substantially outperform humans, preventing them to be trivially used as substitutes of the latter, and, that, moreover, the more the parameters of the LLM the better its performance.

The rest of the paper is organized as follows. In Section \ref{sec:background}, we review the Peira framework. In Section \ref{sec:study}, we describe the experimental study. In Section \ref{sec:related} we offer a brief survey of related work and conclude in Section \ref{sec:conclusions} with future directions.

